\chapter{결론}

PRISM-Hub는 800 V급 차량의 충전 호환과 다중 전압 보조전원을 하나의 직류 전력 구조로 다루며, 최종 구조는 \textbf{개선 구조 R2}(Stage A 중심, 조정형 Stage C, 팩 등급 조건부 창 조정 모듈)가 되었다. 연구 질문의 답은 다음과 같다.

\begin{itemize}
\item \textbf{Q1(정격 분리): 긍정적.} 기존 전 범위 조정 구조에서 조정단 링크전력을 8 kW 이내로 유지하며 충전전력을 주 경로로 전달할 수 있음을 확인했다. R2는 조정이 필요한 팩 등급에만 창 조정 모듈을 두며, 충전기가 V\textsubscript{A} 노드에 붙으므로 충전 중 Stage A 정지를 위한 V\textsubscript{A} 클램프와 정지 순서를 함께 정했다(12장).
\item \textbf{Q2(강건 제어와 부하 관리): 부분적으로 긍정적(H2 부분 충족, 조건: 새 인터페이스 요구).} 같은 제어기가 Tier-1의 공차·지연 기준을 지켰지만, 640 V 복합 공차 Tier-2는 평균모델 기준 93.75 A에서 102/128(스위칭 모델 100 A 기준 128/128)이고, 같은 조합에 +3 µs를 준 스위칭 진단도 114/128이었다(E46). 상태기계는 Stage B 단일 고장에서 Stage A를 유지하며 우회하고 이중 고장에서는 Stage A를 정지하며, 지정한 차단 가능 고장과 메시지 지연·손실 74/74에서 강등 순서가 성립했다(범위는 18.4절). 최종 클램프 회로(TB33)의 오동작 사례 C33은 평균모델과 달랐다(불일치로 기록, 12.6절). 48 V 회생 차단, Stage C 일시 정지, 이중 고장 400 V 부하 정지(3 ms)는 새 요구로 실제 응답은 차량 수준에서 확인해야 하며, 일시 정지 동안 48 V는 배터리 전압을 따르므로(±2 \% 조정 미유지, 10.4절) 48 V 연속성은 배터리를 넣은 결합 시험의 과제다.
\item \textbf{Q3(기준 구조 대비 이점): 원래 구조는 부정적, 개선 구조 R2는 제한적.}
\begin{itemize}
\item \textbf{원래 구조(③): 부정적.} 전 범위 400 V 조정은 손실·스위치 VA·스위치 위치 모두에서 조정 없는 구조보다 컸다(H5 기각, ③ 144.5 대 ④ 101.5 W). 원인은 Stage B의 경부하 고정 손실이다.
\item \textbf{개선 구조 R2: 충전 경로 손실에서만 이점이 있다.} 375 A 충전 손실은 전용 승압기(16–50 kHz)의 0.32–0.51배, 모터 승압 기준점의 0.50–0.51배이고, 상저항 하한과는 같은 수준이다(0.92–1.02배, 920 V에서 12 W 큼). ⑥과의 총손실 교점은 800 V·375 A 기준 주행 1시간당 400–500 V급 충전량으로 모터 승압 기준점(상저항 10 mΩ)에서 4.91 kWh, 상저항 ≥ 5 mΩ 격자 최소에서 177.51 kWh다(920 V 격자 최소와는 양의 교점 없음, 15.2절).
\item \textbf{보조 경로 손실은 이점이 없다.} 400 V급 보조부품을 쓰는 기존 방식(⑤·①)과 같은 수준이며(공정 기준 −1.4 \%, 손실 모델 계수에 따라 −26.5 … +10.0 \%; E\textsubscript{oss} 정정 뒤 ⑤가 명목상 약간 작고, 0.5 kW 운전시간 0.94 \%p 이동으로 순위가 바뀜, E46), 800 V급 보조부품 설계(⑥)보다 4–33 W 크다. 창 조정 모듈도 같은 칸의 다른 성립 구조와 같은 수준이다(7장, E42).
\item \textbf{전력변환 스위치 위치는 더 많다}(④ 42 대 ① 24, ⑤ 12, ⑥ 20; 비상 클램프·구동 인버터 제외, 바이패스 포함). 덧붙여 영전류 게이트의 150·100 kHz 후보는 당시 기준에서 0/30이었으나 25 kHz의 유일성을 보인 것은 아니다(13장, 18.4절).
\end{itemize}
\end{itemize}

2:1 단의 충전·보조부하 겸용 개념은 선행이 있으므로 기여는 구조 개념이 아니라 성립 조건의 정량화, 같은 노드 배치의 회로 검증, 팩 등급 조건부 창 조정 모듈과 고장 강등, 설계 규칙, 부정적 결과(H5 기각)다(15장). Stage A 손실 모델 오류로 설계를 25 kHz로 바꿔 이전 결론을 다시 계산했고(13·14장), 불통과는 지우지 않았다. 모든 판정은 해석·시뮬레이션·데이터시트 대조 범위의 결과이며, 하드웨어 실증과 새 인터페이스 요구·조건부 및 단서 부품 확인은 다음 단계 과제다(19장).
